% Feedback loop diagram for Chapter 2 - Ideal Assumptions
\begin{figure}[htbp]
\centering
\begin{tikzpicture}[node distance=2cm]
    % Nodes
    \node[input] (input) {};
    \node[sum, right=1.2cm of input, fill=white] (sum) {};
    \node at (sum) {\small $\Sigma$};
    \node[controllerblock, right=1.5cm of sum] (controller) {Controller};
    \node[block, right=1.8cm of controller, minimum width=5em] (actuator) {Actuator};
    \node[plantblock, right=1.8cm of actuator] (plant) {Plant};
    \coordinate[right=1.5cm of plant] (output);
    \node[sensorblock, below=1.2cm of plant] (sensor) {Sensor};

    % Signal paths
    \draw[signal] (input) -- node[above] {$r(t)$} node[below, pos=0.85] {\small $+$} (sum);
    \draw[signal] (sum) -- node[above] {$e(t)$} (controller);
    \draw[signal] (controller) -- node[above] {$u(t)$} (actuator);
    \draw[signal] (actuator) -- (plant);
    \draw[signal] (plant) -- node[above] {$y(t)$} (output);

    % Feedback path
    \coordinate[right=0.8cm of plant] (fb1);
    \node[branch] at (fb1) {};
    \draw[thick] (fb1) |- (sensor);
    \draw[signal] (sensor) -| node[left, pos=0.9] {\small $-$} (sum);

    % Ideal assumption boxes
    \node[above=0.3cm of actuator, font=\footnotesize\itshape, text=UofSCHorseshoe] {Ideal: $G_a(s) = 1$};
    \node[below=0.3cm of sensor, font=\footnotesize\itshape, text=UofSCHorseshoe] {Ideal: $H(s) = 1$};
\end{tikzpicture}
\caption{Standard feedback control loop showing the plant, controller, actuator, and sensor. Under ideal assumptions, the actuator and sensor transfer functions are unity.}
\label{fig:feedback_loop_ideal}
\end{figure}
